\documentclass[10pt,a4paper]{amsart}
\usepackage[margin=19mm]{geometry}
\usepackage{amsmath,amssymb,mathtools}
\usepackage[scr=rsfs,cal=euler]{mathalfa}
\usepackage{xfrac}
\usepackage[unicode,hidelinks,bookmarks=false]{hyperref}
\newcommand{\ie}{i.e.\ }
\newcommand{\cf}{cf.\ }
\newcommand{\resp}{respectively\ }
\newcommand{\Subsection}{\subsection}
\providecommand{\nobreakdash}{\nobreak\hbox{-}}
\setlength{\emergencystretch}{2em}
\multlinegap0pt
\usepackage{bm}
\usepackage{bbm}
\usepackage{dsfont}
\usepackage{xspace}
\usepackage{cancel}
\usepackage{mathtools}
\usepackage{amsthm}
\makeatletter
\@ifundefined{thm}{\newtheorem{thm}{Theorem}}{}
\makeatother
\title[The Lax-Wendroff theorem for Patankar-type methods applied t]{The Lax-Wendroff theorem for Patankar-type methods applied to hyperbolic conservation laws}
\author{Janina Bender, Thomas Izgin, Philipp \"Offner, Davide Torlo}
\date{Source: arXiv:2505.08387v2 (CC BY 4.0)}
\begin{document}
\maketitle

\noindent\emph{Source and licence.} The Lax-Wendroff theorem for Patankar-type methods applied to hyperbolic conservation laws. Version arXiv:2505.08387v2 (CC BY 4.0), distributed under the Creative Commons Attribution 4.0 International licence, as recorded in the arXiv licence metadata for this exact version and shown on the abstract page https://arxiv.org/abs/2505.08387v2. Full rights notice: licenses/arxiv-2505.08387-CC-BY-4.0.txt.\par\smallskip

\noindent\textit{Editorial scope.} Counted unit: the Thm whose source label is \texttt{eq:implicit\_CFL} in the article (source0.tex lines 1371--1381, source label \texttt{eq:implicit\_CFL}), delivered together with the article's own complete original proof of it (source0.tex lines 1382--1418, one \texttt{proof} environment of 37 lines). No mathematics is written, repaired or re-derived: statement and proof are the article's own text line for line. Required context retained uncounted: none. Every definition, display and label of those lines is retained unchanged. Omitted from this package and not redistributed: 1--1370, 1419--1457. Nothing inside a retained excerpt is omitted or altered: every retained body is delivered exactly as the article's lines are written, so no omission receipt of any kind accompanies this package. Citations: the retained excerpts carry no citation command at all, so no bibliography entry of the article is transcribed and no reference is redistributed. Reference adaptation: none was needed and none was made; every reference inside the delivered unit resolves inside it, so no unresolved reference or citation is produced. Printed slips kept literal: no printed word, symbol, hypothesis or number of the counted statement or of its proof was altered. Layout and class: the article's own class is \texttt{elsarticle} with packages including geometry, framed, multirow, csquotes, amsmath, amssymb, commath, mathtools, bbm, amsthm, relsize, adjustbox, algorithm, algpseudocode; the wrapper is the lane's pinned \texttt{amsart} editorial wrapper at 10pt on A4, which restates the theorem-like environments the retained text needs and adds the article's own macro definitions that the retained text needs (none), with extra wrapper packages (bm, bbm, dsfont, xspace, cancel, mathtools). The delivered numbering is the unit's own. Assets: no retained span contains a figure environment, a graphics-inclusion command, a PDF-page inclusion, a table or a TikZ picture; no image file is copied into this package, the package declares no asset dependency and no image file is redistributed with it. Licence: version arXiv:2505.08387v2 of this article is distributed under the Creative Commons Attribution 4.0 International licence, as recorded in the arXiv licence metadata for this exact version and shown on the abstract page https://arxiv.org/abs/2505.08387v2. A full rights notice is delivered at licenses/arxiv-2505.08387-CC-BY-4.0.txt. Characters: every retained excerpt is pure ASCII, so the delivered engine, XeLaTeX with its default Latin Modern text font, reports no missing character.
\begin{thm}
	Consider the upwind method with the following PDS distribution
	\begin{equation}
		\partial_t U_i = - \frac{U_i^2-U_{i-1}^2}{2\Delta x}, \qquad \text{with } p_{i,i-1}=\frac{U_{i-1}^2}{2\Delta x}, \quad d_{i,i+1}=\frac{U_{i}^2}{2\Delta x}.
	\end{equation}
	The MPE scheme applied to this semidiscretization is TVD with the \textit{implicit }CFL condition
    \begin{equation}\label{eq:implicit_CFL}
        \frac{\Delta t^n}{\Delta x} U_i^{n+1} \leq 2 \qquad \forall i,
    \end{equation}
    for all times $n$ with $\Delta t ^n=t_{n+1}-t_n$.
\end{thm}

\begin{proof}
	Let us first write down explicitly the MPE scheme applied to the previous semidiscretization:
	\begin{align}
		U^{n+1}_i &= U^n_i - \frac{\Delta t}{2\Delta x}\left( (U_i^n)^2 \frac{U_i^{n+1}}{U_i^n}-(U_{i-1}^n)^2 \frac{U_{i-1}^{n+1}}{U_{i-1}^n} \right)\\
		&=U^n_i - \frac{\Delta t}{2\Delta x}\left( U_i^n U_i^{n+1}- U_{i-1}^n U_{i-1}^{n+1}\right) =U^n_i - \frac{\Delta t}{2\Delta x}\left( U_i^n (U_i^{n+1}-U_{i-1}^{n+1})+ (U^n_i-U_{i-1}^n) U_{i-1}^{n+1}\right) .
	\end{align}
	Now, applying the difference of two consecutive values, we obtain
	\begin{align}
    \begin{split}
	&U^{n+1}_i-U^{n+1}_{i-1} =U^n_i -U^n_{i-1}\\
        &- \frac{\Delta t}{2\Delta x}\left( U_i^n (U_i^{n+1}-U_{i-1}^{n+1})+ (U^n_i-U_{i-1}^n) U_{i-1}^{n+1}  
		-U_{i-1}^n (U_{i-1}^{n+1}-U_{i-2}^{n+1})- (U^n_{i-1}-U_{i-2}^n) U_{i-2}^{n+1}\right)
    \end{split}
	\end{align}
	and bringing on the left hand side the implicit term that we can collect, we obtain
	\begin{align}
    \begin{split}
		&(U^{n+1}_i-U^{n+1}_{i-1})\left(1+\frac{\Delta t}{2\Delta x}U^n_i\right) \\
        =&(U^n_i -U^n_{i-1})\left(1-\frac{\Delta t}{2\Delta x}U^{n+1}_{i-1}\right)+ \frac{\Delta t}{2\Delta x}\left( 
		U_{i-1}^n (U_{i-1}^{n+1}-U_{i-2}^{n+1})+(U^n_{i-1}-U_{i-2}^n) U_{i-2}^{n+1}\right).
    \end{split}
	\end{align}
	Now, using \eqref{eq:implicit_CFL} and the positivity of the solution, 
	we obtain	
	\begin{align}
		\begin{split}
			\sum_i |U^{n+1}_i-U^{n+1}_{i-1}|&\left(1+\frac{\Delta t}{2\Delta x}U^n_i\right) \leq\sum_i|U^n_i -U^n_{i-1}|\left(1-\frac{\Delta t}{2\Delta x}U^{n+1}_{i-1}\right)\\
			&+ \sum_i\frac{\Delta t}{2\Delta x}
			U_{i-1}^n |U_{i-1}^{n+1}-U_{i-2}^{n+1}|+ \sum_i\frac{\Delta t}{2\Delta x}|U^n_{i-1}-U_{i-2}^n|U_{i-2}^{n+1}
		\end{split}\\
		\begin{split}
			\sum_i |U^{n+1}_i-U^{n+1}_{i-1}| &\left(1+\frac{\Delta t}{2\Delta x}U^n_i -\frac{\Delta t}{2\Delta x}
			U_{i}^n \right)  \leq\sum_i|U^n_i -U^n_{i-1}|\left(1-\frac{\Delta t}{2\Delta x}U^{n+1}_{i-1}+ \frac{\Delta t}{2\Delta x}U_{i-1}^{n+1} \right)
		\end{split}\\
		\Longleftrightarrow TV(U^{n+1})&=\sum_i |U^{n+1}_i-U^{n+1}_{i-1}|\leq \sum_i|U^n_i -U^n_{i-1}|=TV(U^{n}).
	\end{align}
\end{proof}

\end{document}
